\documentclass[10pt]{article}
\usepackage[T1]{fontenc}
\usepackage[utf8]{inputenc}
\usepackage[french]{babel}
\usepackage{amsmath,amssymb,mathrsfs,mathpazo}
\usepackage{geometry}
\geometry{a4paper,margin=2cm}
\usepackage{graphicx,array,tabularx,tabularray,colortbl,multirow,multicol,booktabs,enumitem,pifont,ulem,fancybox,diagbox,cancel}
\usepackage{pstricks,pstricks-add,pst-plot,pst-tree,pst-node,pst-text,pst-func,pst-math,pst-bspline,pst-3dplot,pst-eucl,pst-solides3d}
\let\pscalculate\undefined
\usepackage{pst-bezier}
\usepackage{tikz,pgfplots,tkz-tab}
\usetikzlibrary{arrows,arrows.meta,patterns,calc,positioning,trees,decorations.markings,shapes.geometric,shapes.misc,shapes.arrows,intersections,decorations.pathmorphing,decorations.pathreplacing,backgrounds,shadows,babel}
\pgfplotsset{compat=1.18}
\usepgfplotslibrary{fillbetween}
\usepackage[np]{numprint}
\usepackage{listings,makecell,pas-tableur,twoopt}
\let\barre\undefined
\usepackage[pstricks]{tablvar}
\definecolor{darkgreen}{RGB}{0,100,0}
\providecommand{\red}{\color{red}}
\providecommand{\blue}{\color{blue}}
\providecommand{\green}{\color{green}}
\providecommand{\black}{\color{black}}
\providecommand{\white}{\color{white}}
\providecommand{\gray}{\color{gray}}
\providecommand{\cyan}{\color{cyan}}
\providecommand{\magenta}{\color{magenta}}
\providecommand{\yellow}{\color{yellow}}
\providecommand{\violet}{\color{violet}}
\providecommand{\orange}{\color{orange}}
\providecommand{\euro}{\texteuro}
\providecommand{\EUR}{\texteuro}
\providecommand{\ieme}{\textsuperscript{e}}
\providecommand{\degres}{\ensuremath{^\circ}}
\providecommand{\degree}{\ensuremath{^\circ}}
\providecommand{\vv}[1]{\overrightarrow{#1}}
\providecommand{\Coord}[1]{#1}
\providecommand{\decosix}{}
\providecommand{\starredbullet}{\ensuremath{\bullet}}
\pagestyle{empty}
\setlength{\parindent}{0pt}
\usepackage{adjustbox}\DeclareUnicodeCharacter{2236}{:}\DeclareUnicodeCharacter{2190}{\ensuremath{\leftarrow}}
\providecommand{\R}{\mathbb{R}}
\providecommand{\N}{\mathbb{N}}
\providecommand{\Z}{\mathbb{Z}}
\providecommand{\Q}{\mathbb{Q}}
\providecommand{\D}{\mathbb{D}}
\providecommand{\C}{\mathbb{C}}
\providecommand{\e}{{\mathrm{e}}}
\providecommand{\vect}[1]{\overrightarrow{#1}}
\providecommand{\vectt}[1]{\overrightarrow{\mathrm{#1}}}
\providecommand{\vvt}[1]{\overrightarrow{\mathrm{#1}}}
\providecommand{\np}[2][]{{#2~#1}}
\providecommand{\vv}[1]{{\overrightarrow{#1}}}
\renewcommand{\d}{{\mathrm{d}}}
\providecommand{\degres}{{^\circ}}
\providecommand{\degre}{{^\circ}}
\providecommand{\blue}{\color{blue}}
\providecommand{\red}{\color{red}}
\providecommand{\enskip}{\quad}
\providecommand{\llbracket}{\lbrack\!\lbrack}
\providecommand{\rrbracket}{\rbrack\!\rbrack}
\providecommand{\euro}{{\text{€}}}
\providecommand{\no}{{\text{n°}}}
\providecommand{\interval}[2]{{[#1;#2]}}
\providecommand{\intervFO}[2]{{[#1;#2[}}
\providecommand{\dotfill}{\ldots}

\providecommand{\ds}{\displaystyle}
\providecommand{\textding}[1]{#1}
\providecommand{\Oij}{\(\mathrm{O};\vec{\imath},\vec{\jmath}\)}
\providecommand{\Oijk}{\(\mathrm{O};\vec{\imath},\vec{\jmath},\vec{k}\)}
\providecommand{\Ouv}{\(\mathrm{O};\vec{u},\vec{v}\)}
\providecommand{\intervFF}[2]{\ensuremath{[#1;#2]}}
\renewcommand{\euro}{\ifmmode\text{\texteuro}\else\texteuro\fi}
\DeclareUnicodeCharacter{00B0}{\ifmmode^{\circ}\else\textdegree\fi}
\DeclareUnicodeCharacter{20AC}{\ifmmode\text{\texteuro}\else\texteuro\fi}
% Énoncés : Éducation nationale. Transcription et corrigés : APMEP, Denis Vergès et contributeurs. Découpage et classement : CoopMaths, Sébastien Lozano, Éric Elter et contributeurs. Adaptation de la présentation : Maths Explorer.
% CC BY-SA : https://coopmaths.fr/www/about/
% https://coopmaths.fr/alea/static/bac/2024/tex/bac_2024_06_sujet2_polynesie_3_cor.tex
\begin{document}
\noindent\textit{APMEP / CoopMaths. CC BY-SA. Présentation adaptée par Maths Explorer.}\par\medskip
\medskip

On considère la suite $\left(u_n\right)$ définie par : 

\begin{center} $u_0 = 8$ \:  et pour tout entier naturel \:\:$n,\:\:~u_{n+1} = u_n -\ln\left(\dfrac{u_n}{4}\right)$.\end{center}

\begin{enumerate}
\item 
	\begin{enumerate}
		\item ~
		
$\bullet~u_1 = 8 - \ln 2 \approx 7,31$ ;
		
$\bullet~u_2 = 8 - \ln 2  \ln {8 - \ln 2}{4}- \approx 6,70$.

		\item 












	

\texttt{mystere(10)} donne la somme des premiers termes de la suite de $u_0$ à $u_9$, soit 

$u_0 + u_1 + \ldots + u_9 = \displaystyle\sum_{k=0}^n u_k$.
		\item 
Il suffit en ligne 7 de remplacer S par (S /k)
	\end{enumerate}
\item 

\[ f(x) = x - \ln\left(\dfrac{x}{4}\right).\]













	

	

La fonction $f$ somme de fonctions dérivables sur $[0~;~+\infty[$ est dérivable et sur cet intervalle en posant $u(x) = \dfrac x4$, d'où $u'(x) = \dfrac14$,

$f'(x) = 1 - \dfrac14 \times \dfrac{1}{\dfrac x4} = 1 - \dfrac14 \times 4 \dfrac 1x  = 1 - \dfrac 1x  = \dfrac{x - 1}{x}$

Comme $x > 0$, le signe de $f'(x)$ est celui de $x - 1$ :

$\bullet~$ si $0< x < 1$, alors $x - 1 < 0$  : la fonction $f$ est décroissante sur $]0~;~1[$ ;

$\bullet~$ si $x > 1$, alors $x - 1 > 0$ : la fonction $f$ est croissante sur $-1~;~+ \infty[$

$\bullet~$si $x = 1$, alors $f'(1) = 0$ et $f(1) = 1 - \ln \dfrac14 = 1 - (-\ln 4) = 1 + \ln 4$ est le minimum de la fonction $f$ sur $[0~;~+ \infty[$.
\end{enumerate}

Dans la suite de l'exercice, on remarquera que pour tout entier naturel $n$,\: $u_{n+1} = f\left(u_n\right)$.

\begin{enumerate}[resume]
\item 
	\begin{enumerate}
		\item 


\emph{Initialisation} : on a vu que $u_1 \approx 7,3$ et comme $u_0 = 8$, on a donc :

\[1 \leqslant u_{1} \leqslant u_0\]
l'encadrement est vrai au rang 0.

\emph{Hérédité} : soit $n \in \N$, tel que  $1 \leqslant u_{n+1} \leqslant u_n$.

On sait que sur l'intervalle $[1~;~+ \infty[$ la fonction $f$ est croissante, donc on a 

$f(1) \leqslant f\left(u_{n+1}\right) \leqslant f\left(u_{n}\right)$ .

Or $f(1) = 1 - \ln \dfrac14 \approx 2,39$, donc $f(1) \geqslant 1$ et on a 

$1 \leqslant u_{n+2} \leqslant u_{n+1}$ : l'encadrement est vrai au rang $n + 1$.

\emph{Conclusion} : l'encadrement est vrai au rang 0 et s'il est vrai au rang $n \in \N$ il l'est encore au rang suivant : d'après le principe de récurrence :

pour tout entier naturel $n$, on a : $1 \leqslant u_{n+1} \leqslant u_n.$

	\item L'encadrement précédent montre que : pour tout naturel $n$

$\bullet~ u_{n+1} \leqslant u_n$ : la suite $\left(u_n\right)$ est décroissante ;

$\bullet~ 1 \leqslant u_n$ :  la suite $\left(u_n\right)$ est minorée par 1.

On sait qu'alors la suite $\left(u_n\right)$ converge vers une limite $\ell \geqslant 1$




		\item 
$f(x) = x \iff x - \ln \left(\dfrac x4 \right) = x \iff 0 = \ln \left(\dfrac x4 \right) \iff \dfrac x4  = 1 \iff x = 4$.
		\item 
Par continuité de la fonction $f$ (car elle est dérivable sur $]0~;~+ \infty[$), la relation 

$u_{n+1} = u_n - \ln \left(\dfrac{u_n}{4}\right)$, donne puisque $\displaystyle\lim_{ n \to + \infty} u_n = \ell$ :

$\ell = \ell - \ln \left(\dfrac{\ell}{4}\right)$ : d'après la question précédente $\ell = 4$.

	\end{enumerate}
\end{enumerate}

\bigskip

\end{document}
